% Original diagram showing the structural Markov blanket in an undirected graph.
\begin{tikzpicture}[
  x=1mm,y=1mm,
  center/.style={circle,draw=BookAmber,fill=BookAmber!18!white,double,
    double distance=0.8pt,minimum size=10mm,inner sep=0pt,font=\small,
    line width=0.8pt},
  neighbor/.style={circle,draw=BookTeal,fill=FigureModel!10,
    minimum size=8mm,inner sep=0pt,font=\small,line width=0.9pt},
  remote/.style={rectangle,draw=BookMuted,fill=BookPaper,
    minimum width=8mm,minimum height=8mm,inner sep=0pt,font=\small,
    line width=0.9pt},
  graph edge/.style={draw=BookInk,line width=0.9pt},
  heading/.style={font=\heiti\footnotesize,text=BookInk},
  note/.style={font=\scriptsize,text=BookMuted,align=center},
  every node/.style={text=BookInk}
]
  \node[heading] at (10,52) {中心节点};
  \node[heading] at (47,52) {邻居集合 $N(i)$};
  \node[heading] at (91,52) {远端非邻居};

  \node[center] (i) at (10,25) {$i$};

  \node[neighbor] (n1) at (47,39) {$n_1$};
  \node[neighbor] (n2) at (47,25) {$n_2$};
  \node[neighbor] (n3) at (47,11) {$n_3$};

  \node[remote] (r1) at (91,45) {$r_1$};
  \node[remote] (r2) at (91,32) {$r_2$};
  \node[remote] (r3) at (91,18) {$r_3$};
  \node[remote] (r4) at (91,5) {$r_4$};

  \draw[graph edge] (i)--(n1);
  \draw[graph edge] (i)--(n2);
  \draw[graph edge] (i)--(n3);

  \draw[graph edge] (n1)--(r1);
  \draw[graph edge] (n1)--(r2);
  \draw[graph edge] (n2)--(r2);
  \draw[graph edge] (n2)--(r3);
  \draw[graph edge] (n3)--(r3);
  \draw[graph edge] (n3)--(r4);

  \node[note] at (27,-1) {离开 $i$ 的第一步必在 $N(i)$ 中};
\end{tikzpicture}
